\documentclass[11pt,a4paper]{article}
\usepackage[margin=25mm]{geometry}\usepackage{amsmath,amssymb,amsthm}
\newtheorem{lemma}{Lemma}\newtheorem{theorem}{Theorem}
\title{A product obstruction to chiral quotients}
\author{Human Compression project}\date{}
\begin{document}\maketitle
We use the previously established normalizer description of
colour-preserving automorphisms of a connected coset flag graph.
The following elementary consequence removes an entire class of
regular covers before any subgroup enumeration. No priority claim is
made for the observation.

Let a regular polytope have string C-group
$W=\langle r_0,\ldots,r_{d-1}\rangle$. Suppose a quotient polytope
has flags $Nw$, with adjacency $Nw\mapsto Nwr_i$. The assumption that
these are its actual flags is essential.

\begin{lemma}
Suppose the colours partition into nonempty sets $I,J$, every generator
in $I$ commutes with every generator in $J$, and
$W=W_I\times W_J$. Then
\[
 N=(N\cap W_I)\times(N\cap W_J).
\]
\end{lemma}
\begin{proof}
Write $n=ab\in N$ with $a\in W_I$ and $b\in W_J$. The flag
$Na=Nb^{-1}$ is reachable from $N$ both by an $I$-coloured path and
by a $J$-coloured path. The first path preserves every face of rank in
$J$; the second preserves every face of rank in $I$. Hence this flag
contains precisely the faces of the flag $N$, so it is $N$ itself.
Thus $a,b\in N$. Every element of $N$ therefore lies in the product
of the two displayed intersections. The reverse inclusion is automatic,
and the factors commute and meet trivially.
\end{proof}

\begin{theorem}
A regular cover with an adjacent product of order two has no chiral
polytope quotient with the stated coset flag action.
\end{theorem}
\begin{proof}
Suppose $r_jr_{j+1}$ has order two. Since the generators are
involutions, they commute. The string relations imply that every
generator in $I=\{0,\ldots,j\}$ commutes with every generator in
$J=\{j+1,\ldots,d-1\}$. The intersection property gives
$W_I\cap W_J=1$, and generation gives $W=W_I\times W_J$.
By the lemma, $N=N_I\times N_J$, where $N_I=N\cap W_I$ and
$N_J=N\cap W_J$. Consequently
\[
 N_W(N)=N_{W_I}(N_I)\times N_{W_J}(N_J).
\]

If the quotient were chiral, its automorphism group would have exactly
two flag orbits, and every adjacency would exchange them. Any product
$r_ir_k$ of two generators therefore returns the base flag to its
automorphism orbit, so $r_ir_k\in N_W(N)$ by the normalizer theorem.
Take $i\in I$ and $k\in J$. The displayed product decomposition of
the normalizer implies $r_i\in N_{W_I}(N_I)$ and
$r_k\in N_{W_J}(N_J)$. Varying $i$ and $k$ shows that all generators
normalize $N$. Thus $N_W(N)=W$, and the quotient has one flag orbit,
a contradiction.
\end{proof}

The same proof applies to any partition satisfying the lemma, whether
or not it arises from a single break in a string. It does not assert
that an arbitrary subgroup of a direct product splits; the flag
interpretation is what forces that conclusion. Nor does it identify
an arbitrary orbit poset with its coset graph.

For ranks one and two there is no separate chiral case. A connected
rank-two flag graph is an alternating cycle; rotations by two edges
and a reflection preserving colours act transitively on its vertices.
Rank one has a single edge. Thus the remaining small-cover search may
restrict to rank at least three and to Schlaefli entries all at least
three, while still checking the polytopal quotient condition.

This argument uses the source's regular-cover framework and the
previous normalizer proof. It is a general exclusion, not a completed
atlas census or a solution of a new open question. Completeness of the
remaining finite cover roster remains an explicit input obligation.
\end{document}
